\label{chap:conclusion}

There are several limitations to be addressed in the BioseqDB:
\begin{itemize}
    \item Grouped-MPF does not correctly work for update on insert.
    \item Grouped-MPF uses woraround with writing a file to efficiently support joins.
    \item The data is fetched only in one Postgres process, which causes the one core of CPU to be bottlenecked.
    \item Some kernels contains probable mistakes that slows down  processing for any cases than k=15,w=10 significantly.
    \item The pipeline code is overcomplicated and current type-driven template system causes the build to last up to 2 hours for optimsied builds on some machines.
    \item The storage is optimised only for large sequences, so short sequences overuse memory significantly.
    \item Buggy conditions cause minimizers to be recomputed instead of being loaded from the index.
    \item Index format contains too many fields of which most are not used and the loading code heavily relies of bitwise operations.
    \item Cache is experimental and didn't go proper testing, so I have no confirmation if it works reliably.
    \item Benchmarking was done only on synthetic datasets instead of sequences from NCBI, which can mean that practical speedups are significantly lower.
    \item Currrent join logic assumes reuse of target fragments, which means that BioseqDB can handle large table vs. small table, but will be slow for multi-alignment (large table vs. large table)
    \item The error messages created by BioseqDB are hard to interpret and not very helpful to the user
    \item Compilation requires expiermental version of tools e.g latest CMake version with special flag enabled to build C++ std modules.
\end{itemize}

However, the project can be classified as a success:
\begin{itemize}
    \item BioseqDB is first in literature aim to integrate existing database software with genomic toolset.
    \item BioseqDB supports abitratly large sequences of any sizes, so it doesn't need any splitting step 
    \item Most of the common operations like slicing, ordering, transforming sequences can be done in SQL instead of external tools
    \item Most of the algorithms were ported by me based on their original source code, so now they're supported on any architecture.
    \item BioseqDB aligner pipeline can be called as standalone binary directly.
    \item I provided extensive analysis of algoritms that are poorly described in literature, so the provided illustrations and code listings can be used as a base to develop new software more easily and great learning material.
    \item The overall performance was very satisfactory and speedups of x2, x10 to x100 on heavy SWAPped systems were achieved.
\end{itemize}

\section{Acknowledgments}

I would like to thank to my supervisor prof. dr hab. Krzysztof Stencel, whos expertise in database systems was an inspiration for me to execute this project.
I dedicate this work to Wiesław Styczyński and Agnieszka Styczyńska, who supported me with good words in times of doubts.

\section{Funding}

The work was self-funded and total spending was around $1000\$$ for AWS infrastructure and hardware rentals.